%% LyX 2.4.4 created this file.  For more info, see https://www.lyx.org/.
%% Do not edit unless you really know what you are doing.
\documentclass[english]{article}
\usepackage[T1]{fontenc}
\usepackage[latin9]{inputenc}
\usepackage{float}
\usepackage{enumitem}
\usepackage{amsmath}
\usepackage{amsthm}
\usepackage{graphicx}
\usepackage{geometry}
\geometry{verbose,tmargin=1in,bmargin=1in,lmargin=1in,rmargin=1in}
\PassOptionsToPackage{normalem}{ulem}
\usepackage{ulem}

\makeatletter
%%%%%%%%%%%%%%%%%%%%%%%%%%%%%% Textclass specific LaTeX commands.
\numberwithin{equation}{section}
\numberwithin{figure}{section}
\newlength{\lyxlabelwidth}      % auxiliary length 

%%%%%%%%%%%%%%%%%%%%%%%%%%%%%% User specified LaTeX commands.
\usepackage{fancyhdr}
\usepackage{lastpage}
\pagestyle{fancy}
\usepackage{enumitem}
\usepackage{ifthen}
%sets math fonts to be more readable
\everymath=\expandafter{\the\everymath\displaystyle}
%\DeclareMathSizes{10}{10}{10}{10}
\usepackage{multicol}
% Added by lyx2lyx
\setlength{\parskip}{\smallskipamount}
\setlength{\parindent}{0pt}

\makeatother

\usepackage{babel}
\begin{document}
\renewcommand{\author}{Dr.\,James Ashe}

\newcommand{\classNum}{232}

\newcommand{\classSec}{C}

\newcommand{\nameType}{Name:}

\newcommand{\examType}{Exam 3}

\renewcommand{\date}{\today} %%\today

%%First page's header%%%%%%%%%%%%%%%%%%%%%%
\fancypagestyle{firststyle} %%header settings for first pagr
{
	\fancyhead[L]{MTH \classNum{}(\classSec{}) \author{}\\
	\textbf{\examType{}} \date{}}
	\fancyhead[C]{\textbf{\nameType}}
	\setlength{\headsep}{14pt}
}
%%%%%%%%%%%%%%%%%%%%%%%%%%%%%%%%%%%%%%%%%%

%%Next page's header%%%%%%%%%%%%%%%%%%%%%%%
\setlist{nolistsep}
\lhead{\classNum{}(\classSec{})} 
\chead{\textbf{\nameType}} 
\cfoot{\thepage{}~of~\pageref{LastPage}} 
\setlength{\headsep}{5pt} 
%%%%%%%%%%%%%%%%%%%%%%%%%%%%%%%%%%%%%%%%%%%

%%Various settings; see the preamble for other settings%%%%%
%%\setenumerate[0]{leftmargin=12pt,itemindent=0pt,labelwidth=0pt}
\renewcommand{\labelenumi}{\textbf{\arabic{enumi}.}}
\renewcommand{\labelenumii}{\textbf{\alph{enumii}.}}
\thispagestyle{firststyle}
%%%%%%%%%%%%%%%%%%%%%%%%%%%%%%%%%%%%%%%%%%

Unless otherwise stated, each problem is worth 4 points. Show all
of your work.

\textbf{\emph{(a) }}\emph{Find the inverse on the specified interval
express it in the form $y=f^{-1}(x)$. }\textbf{\emph{(b)}}\emph{
Verify that $f^{-1}\left(f(x)\right)=x$. }\textbf{\emph{(c)}}\emph{
State the range and domain of $f^{-1}(x)$.}
\begin{enumerate}[resume]
\item $f(x)=\sqrt{x-3}+1$ for $x\geq3$ {[}6 pts{]}\vspace{1.5cm}\vfill
\end{enumerate}
\emph{Find the derivative of the inverse of the following functions
at the specified point on the graph of the inverse function }\emph{\uline{without
finding}}\emph{ $f^{-1}$. }
\begin{enumerate}[resume]
\item \textbf{$f(x)=x^{4}-2x^{3}+1$}, at $(0,1)$ (point on $f^{-1})$.\vfill
\end{enumerate}
\emph{Use the graph of $f(x)$ to sketch $f^{-1}(x)$.}
\begin{enumerate}[resume]
\item {[}6 pts{]} (The dashed line is $y=x$.)\\
\begin{figure}[H]
\centering{}\includegraphics[width=2.5in]{D:/home/jim/JamesAshe/JCSU/Calc_2_MTH_232/images/graph_the_inverse.png}
\end{figure}
\newpage
\end{enumerate}

\emph{Perform the indicated operation.}
\begin{enumerate}[resume]
\item $\frac{d}{dx}\left(\ln x+1\right)$\vfill
\item $\frac{d}{dx}$$3e^{x}\ln x$\vfill
\item $\frac{d}{dx}\left(\ln(\cos x)-2x\right)$\vfill
\item $\frac{d}{dx}$$e^{-2x+x^{2}}$\vfill\newpage
\item $\frac{d}{dt}\frac{e^{t}}{e^{t}+1}$\vfill
\item ${\displaystyle \int\left(3e^{-3x}+\frac{1}{x}\right)\,dx}$\vfill
\item $\frac{d}{dx}5\cdot4^{x}\,dx$\vfill
\item ${\displaystyle \int_{-1}^{1}9^{x}}$ {[}6 pts{]}\vfill\newpage
\item ${\displaystyle \int}xe^{-4x}\,dx$\vfill
\item ${\displaystyle \int x\cos4x\,dx}$\vfill
\item ${\displaystyle \int\frac{4\ln x}{x^{5}}\,dx}$\vfill
\item ${\displaystyle \int_{0}^{\ln1}3xe^{x}\,dx}$ {[}6 pts{]}\vfill
\end{enumerate}

\end{document}
